% Figure: Coulomb counting. Top panel shows the measured discharge
% current (nominal 1.5 A with a load transient near t = 900 s); bottom
% panel shows the running charge removed, the trapezoidal integral of
% the current. Generated to match code/coulomb_counting.py and
% data/coulomb_trace.csv. Time axis in 100 s units for readability.
\begin{figure}[htbp]
\centering
\begin{tikzpicture}

% Top panel: current trace. Nominal 1.5 A plus a Gaussian transient
% of peak ~0.7 A centred at t = 900 s (= 9 in 100 s units), width 40 s.
\begin{axis}[
  name=cur,
  width=11.5cm, height=4.2cm,
  xlabel={},
  ylabel={current $\abs{I}$ (A)},
  xmin=0, xmax=18, ymin=0, ymax=2.6,
  domain=0:18, samples=200,
  axis lines=left, grid=both,
  major grid style={engblack!12}, font=\small,
  xticklabels={},
  title={Measured discharge current},
]
\addplot[engblue, very thick] {1.5 + 0.7*exp(-((x-9)^2)/(2*0.4^2))};
\addplot[engamber, thick, dashed, domain=0:18] {1.5};
\node[font=\scriptsize, anchor=south west, engamber!80!black]
  at (axis cs:0.3,1.55) {nominal $1.5\,$A};
\node[font=\scriptsize, anchor=south, engblue]
  at (axis cs:9,2.25) {load transient};
\end{axis}

% Bottom panel: running charge removed (coulombs). The integral of the
% top panel: linear at slope 150 C per 100 s on the flat regions, with
% an extra rise across the transient. Closed form of the cumulative
% integral (current in A, time in 100 s units, so dt-factor = 100):
%   Q(x) = 150 x + 100 * 0.7 * sqrt(2 pi) * 0.4 * 0.5 *
%          (erf((x-9)/(0.4 sqrt2)) - erf((0-9)/(0.4 sqrt2)))
% The erf term saturates to ~ 70.2 C across the transient. pgfplots has
% no erf, so we approximate the saturating step with a logistic of
% matched midpoint and total rise; visually indistinguishable here.
\begin{axis}[
  at={(cur.south)}, anchor=north, yshift=-10pt,
  width=11.5cm, height=4.6cm,
  xlabel={time $t$ (units of $100\,$s)},
  ylabel={charge removed (C)},
  xmin=0, xmax=18, ymin=0, ymax=2900,
  domain=0:18, samples=200,
  axis lines=left, grid=both,
  major grid style={engblack!12}, font=\small,
  title={Running charge $Q(t)=\int_{0}^{t}\abs{I}\,\dd\tau$
    (trapezoidal)},
]
\addplot[engred, very thick]
  {150*x + 70.2/(1 + exp(-(x-9)/0.42))};
\node[font=\scriptsize, anchor=north west, engred]
  at (axis cs:0.5,2750) {final $\approx 2.77\times 10^{3}\,$C
    $=0.770\,$Ah};
\end{axis}

\end{tikzpicture}
\caption[State of charge by Coulomb counting]{Coulomb counting on a
$1800\,\si{\second}$ cell discharge. The top panel is the measured
current: a nominal $1.5\,\si{\ampere}$ discharge (amber dashed) with a
load transient peaking near $2.2\,\si{\ampere}$ around $t =
900\,\si{\second}$. The bottom panel is the running charge removed, the
trapezoidal integral of the current; its slope is the current, so the
transient appears as a steepening of the otherwise straight
accumulation. The integral reaches $2.77\times 10^{3}\,\si{\coulomb} =
0.770\,\si{\ampere\hour}$, a $25.6\,\%$ drop on a
$3.0\,\si{\ampere\hour}$ cell. The dominant error is not the
$20\,\si{\milli\ampere}$ measurement noise, which the integral
averages down to three parts in ten thousand, but the sensor offset,
which integrates coherently and grows linearly with the integration
window. The trace is reproduced by
\texttt{code/coulomb\_counting.py}.}
\label{fig:vol02:ch06:coulomb-counting}
\end{figure}
